\clearpage
\section{Attention Variants}\label{sec:exp:attention}\label{sec:exp:results:attention}

This section compares jointly trained encoders that differ only in how time and incidence enter attention, which is the comparison for RQ3.
Five configurations are trained with the joint objective of \cref{eq:method:loss_total}.
Emitter-only copies of RoPE-TOA and Bias follow the joint bake-off in \cref{sec:exp:attn_task}.
Hyperparameters follow \cref{sec:exp:hparams}.
Vanilla is the joint encoder of \cref{sec:exp:results:vanilla}.
RoPE-TOA applies rotary encoding on the window-local \gls{toa} in the trunk and in both branches.
RoPE-az is the branch-specific rotary design of \cref{sec:method:attention}.
Despite the name, it does not rotate by azimuth or elevation: the trunk and the mode branch rotate by window-local \gls{toa}, and the deinterleaving branch rotates by the two planar Euclidean incidence components $\tilde{u}^{x},\tilde{u}^{y}$ of \cref{eq:method:aoa_scale}.
Bias subtracts the pairwise physical term of \cref{eq:method:attn_bias,eq:method:attn_bias_trunk} from the attention logits.
RoPE-TOA+Bias uses that same \gls{toa} rotary encoding in the trunk and in both branches, together with the pairwise bias; it does not add the incidence rotation of RoPE-az.

The rows other than Vanilla test whether explicit physical structure in attention improves clustering relative to standard \gls{mha}.
\Cref{tab:exp:attn_em,tab:exp:attn_md} collect the window-averaged scores.
The best value in each column is in bold and ties are bold on every tied row.

\begin{table}[H]
  \centering
  \small
  \caption{Emitter clustering on the test split for jointly trained attention variants.}
  \label{tab:exp:attn_em}
  \begin{tabular}{@{}lccccc@{}}
    \toprule
    Model & $\ARI$ & $\AMI$ & Homogeneity & Completeness & Count error \\
    \midrule
    Vanilla   & $0.948\pm0.176$ & $0.957\pm0.138$ & $0.989\pm0.052$ & $0.951\pm0.166$ & $0.25\pm0.49$ \\
    RoPE-TOA  & $\mathbf{1.000\pm0.003}$ & $\mathbf{0.998\pm0.008}$ & $0.996\pm0.015$ & $\mathbf{1.000\pm0.000}$ & $0.11\pm0.34$ \\
    RoPE-az   & $0.973\pm0.135$ & $0.976\pm0.112$ & $0.994\pm0.042$ & $0.974\pm0.130$ & $0.16\pm0.40$ \\
    Bias      & $\mathbf{1.000\pm0.002}$ & $\mathbf{0.998\pm0.007}$ & $\mathbf{0.997\pm0.013}$ & $\mathbf{1.000\pm0.000}$ & $\mathbf{0.10\pm0.33}$ \\
    RoPE-TOA+Bias & $0.851\pm0.255$ & $0.886\pm0.196$ & $0.984\pm0.067$ & $0.846\pm0.236$ & $0.73\pm0.94$ \\
    \bottomrule
  \end{tabular}
\end{table}

\begin{table}[H]
  \centering
  \small
  \caption{Mode clustering on the test split. Entries follow the same aggregation as \cref{tab:exp:attn_em}.}
  \label{tab:exp:attn_md}
  \begin{tabular}{@{}lccccc@{}}
    \toprule
    Model & $\ARI$ & $\AMI$ & Homogeneity & Completeness & Count error \\
    \midrule
    Vanilla   & $0.985\pm0.107$ & $0.986\pm0.102$ & $0.994\pm0.027$ & $0.988\pm0.103$ & $0.23\pm0.50$ \\
    RoPE-TOA  & $0.991\pm0.075$ & $0.991\pm0.068$ & $0.996\pm0.021$ & $0.992\pm0.071$ & $0.20\pm0.48$ \\
    RoPE-az   & $0.983\pm0.102$ & $0.986\pm0.087$ & $0.989\pm0.047$ & $\mathbf{0.993\pm0.083}$ & $0.24\pm0.51$ \\
    Bias      & $\mathbf{0.992\pm0.081}$ & $\mathbf{0.992\pm0.079}$ & $\mathbf{0.997\pm0.018}$ & $\mathbf{0.993\pm0.079}$ & $\mathbf{0.14\pm0.41}$ \\
    RoPE-TOA+Bias & $0.977\pm0.112$ & $0.981\pm0.101$ & $0.986\pm0.047$ & $0.988\pm0.099$ & $0.32\pm0.60$ \\
    \bottomrule
  \end{tabular}
\end{table}

RoPE-TOA and Bias remove the emitter tail that joint Vanilla leaves in \cref{tab:exp:vanilla,tab:exp:ablation_loss}.
Mean emitter $\ARI$ is $1.000$ to three decimals, $99\%$ of windows sit at or above $0.99$, and none of the windows falls below $0.90$.
Completeness is $1.000$ with no spread, so a true emitter is not split across clusters.
The remaining mistakes are rare undercounts: homogeneity sits just below $1$, and the mean absolute count error is $0.11$ for RoPE-TOA and $0.10$ for Bias.
Emitter-only Vanilla already sits near that ceiling in \cref{tab:exp:ablation_loss} ($\ARI$ $0.999$).
Joint RoPE-TOA and Bias sit slightly higher, with a tighter spread.

RoPE-az improves on Vanilla for emitters ($\ARI$ $0.973$, $4.6\%$ of windows below $0.90$) but does not match TOA-only rotary encoding.
Rotating the deinterleaving branch by planar incidence is not, on this test set, a replacement for rotating every layer by \gls{toa}.
Mode scores occupy a narrow band for Vanilla, RoPE-TOA, RoPE-az, and Bias.
Bias is the tightest there, with the smallest count error.

RoPE-TOA+Bias does not keep the emitter gains of either parent.
Mean emitter $\ARI$ is $0.851$, below joint Vanilla, and $31\%$ of windows fall below $0.90$.
Completeness is $0.846$, so the remaining mistakes are splits.
Mode $\ARI$ stays in the same band as the other joint models.
Stacking the two time mechanisms is not an improvement on this split.

\subsection{Single-Task RoPE-TOA and Bias}\label{sec:exp:attn_task}

The joint scores do not say whether RoPE-TOA and Bias would deinterleave as well without the mode term.
\Cref{tab:exp:attn_em_task} reports emitter $\ARI$ for those encoders trained with $\Loss_{\mathrm{em}}$ only.
Vanilla is the ablation of \cref{tab:exp:ablation_loss}.
RoPE-az and RoPE-TOA+Bias have no emitter-only run.

\begin{table}[H]
  \centering
  \small
  \caption{Emitter $\ARI$ for joint and single-task (emitter-only) Vanilla, RoPE-TOA, and Bias. Joint rows are copied from \cref{tab:exp:attn_em}. Vanilla emitter-only is copied from \cref{tab:exp:ablation_loss}.}
  \label{tab:exp:attn_em_task}
  \begin{tabular}{@{}lcc@{}}
    \toprule
    Model & Joint & Emitter-only \\
    \midrule
    Vanilla   & $0.948\pm0.176$ & $0.999\pm0.012$ \\
    RoPE-TOA  & $1.000\pm0.003$ & $0.998\pm0.036$ \\
    Bias      & $1.000\pm0.002$ & $0.840\pm0.265$ \\
    \bottomrule
  \end{tabular}
\end{table}

Emitter-only RoPE-TOA sits slightly below the joint row.
Joint RoPE-TOA and Bias are at par with emitter-only Vanilla; the margins are small.
The Vanilla joint tax is therefore specific to standard attention.
Emitter-only Bias does not match: mean $\ARI$ falls to $0.840$, below joint Vanilla.
The pairwise bias is not a drop-in deinterleaving-only encoder on this split.

\subsection{Per-Window ARI Distributions}\label{sec:exp:results:hist}

\Cref{fig:exp:hist_em,fig:exp:hist_md} show the empirical CDF of per-window $\ARI$ for Vanilla, RoPE-TOA, RoPE-az, Bias, and RoPE-TOA+Bias, the same five models as \cref{tab:exp:attn_em,tab:exp:attn_md}.
Each curve is the fraction of the $1169$ windows whose $\ARI$ does not exceed the value on the horizontal axis.
The vertical axis is logarithmic so that a tail of a few windows remains visible against the jump at $1$.
$\AMI$, homogeneity, and completeness follow the same right spike and are not repeated as extra pages.

The Vanilla emitter curve leaves the axis well before $0.90$.
RoPE-TOA+Bias leaves earlier still, and $57\%$ of its windows sit at or above $0.99$ against $88\%$ for Vanilla.
RoPE-TOA and Bias stay near zero until just below $1$, and the two curves nearly coincide on this split.
RoPE-az sits between Vanilla and those two shapes.
On the mode branch the five curves rise later and closer together. Bias is the last to leave zero. RoPE-TOA+Bias sits with Vanilla and RoPE-az.
\Cref{fig:app:hist_em_all,fig:app:hist_md_all} add Vanilla single-task on the matching branch.

\begin{figure}[H]
  \centering
  \expfig{hist_em.pdf}
  \caption{Empirical CDF of emitter $\ARI$ over the $1169$ non-overlapping test windows. The vertical axis is logarithmic.}
  \label{fig:exp:hist_em}
\end{figure}

\begin{figure}[H]
  \centering
  \expfig{hist_md.pdf}
  \caption{Empirical CDF of mode $\ARI$ over the same windows as \cref{fig:exp:hist_em}. The vertical axis is logarithmic.}
  \label{fig:exp:hist_md}
\end{figure}

\subsection{Cluster-Count Recovery}\label{sec:exp:results:kcorr}

\Cref{fig:exp:kcorr_em,fig:exp:kcorr_md} are heatmaps of estimated versus true cluster counts for the jointly trained attention variants, with the same reading as \cref{fig:exp:kcorr_ablation_em,fig:exp:kcorr_ablation_md}.

On the emitter branch, Vanilla leaves the diagonal more often than RoPE-TOA, RoPE-az, and Bias, with splits and merges in roughly equal number.
RoPE-TOA never oversplits an emitter window in this set: the off-diagonal mass is only undercounts.
Bias is the same pattern, up to one window.
RoPE-az is closer to the diagonal than Vanilla and still produces some extra clusters.
RoPE-TOA+Bias leaves the diagonal more often than Vanilla, and the extra mass is above it.

The mode heatmaps are diagonal-dominant for every model.
When they miss, they more often undercount than invent labels, and the misses concentrate at the crowded end of the catalogue, where a window can contain a dozen modes (\cref{sec:exp:datasets}).
Bias keeps more mass on the diagonal than Vanilla or RoPE-az.
RoPE-TOA+Bias is noisier there than Bias, with more off-diagonal counts at the high-$M$ end.

\begin{figure}[tbp]
  \centering
  \expfig{kcorr_em.pdf}
  \caption{Estimated versus true number of emitters per test window. Cell values are window counts. The line is $\hat{K}=K$.}
  \label{fig:exp:kcorr_em}
\end{figure}

\begin{figure}[tbp]
  \centering
  \expfig{kcorr_md.pdf}
  \caption{Estimated versus true number of modes per test window. Layout as in \cref{fig:exp:kcorr_em}.}
  \label{fig:exp:kcorr_md}
\end{figure}

\subsection{Model Size and Inference Cost}\label{sec:exp:results:efficiency}

\Cref{tab:exp:efficiency} reports trainable parameters and wall-clock time per test window on the A100 of \cref{sec:exp:setup}.
One window is forwarded at a time after a short GPU warmup, which is excluded from the mean.
Clustering time is \gls{dbscan} on the embeddings of \cref{eq:method:cluster_decoding}.

\begin{table}[H]
  \centering
  \caption{Trainable parameters and mean time per test window. Encoder is the forward pass, or the feature copy for Feature \gls{dbscan}. Clustering is \gls{dbscan} on both branches.}
  \label{tab:exp:efficiency}
  \begin{tabular}{@{}lcccc@{}}
    \toprule
    Model & Parameters & Encoder (ms) & Clustering (ms) & Total (ms) \\
    \midrule
    Feature DBSCAN & $0$ & $0.1$ & $20.9$ & $21.0$ \\
    Vanilla   & $16\,081\,536$ & $1.6$ & $58.2$ & $59.8$ \\
    RoPE-TOA  & $16\,081\,536$ & $2.5$ & $55.9$ & $58.4$ \\
    RoPE-az   & $16\,081\,536$ & $3.1$ & $54.6$ & $57.7$ \\
    Bias      & $16\,081\,536$ & $31.8$ & $56.3$ & $88.1$ \\
    RoPE-TOA+Bias & $16\,081\,536$ & $30.1$ & $52.4$ & $82.5$ \\
    \bottomrule
  \end{tabular}
\end{table}

Rotary encoding adds no learned weights, so both RoPE models match Vanilla.
The physical bias stores one scalar per head, layer, and selected coordinate; those scalars do not change the dumped total at the reported precision.
They do change the encoder time: Bias and RoPE-TOA+Bias spend about $30$~ms in the forward pass, against $1.6$~ms for Vanilla.
Those times include rebuilding the pairwise maps $B^{(p)}_{ij}=\lvert p_i-p_j\rvert$ of \cref{eq:method:attn_bias,eq:method:attn_bias_trunk} in every head of every sub-layer; the matrices depend only on the window coordinates, and no cache was used.

\Gls{dbscan} is about $55$~ms for every neural model and dominates Vanilla and both RoPE variants.
Feature \gls{dbscan} has no encoder weights and clusters in $21$~ms.
Memory for the pairwise maps is discussed in \cref{sec:disc:attention}.

\subsection{Learned Bias Coefficients}\label{sec:exp:results:bias_coef}

Bias and RoPE-TOA+Bias store one scalar per head, layer, and selected coordinate of \cref{eq:method:attn_bias,eq:method:attn_bias_trunk}.
The trunk has only time; each task branch also has an incidence coefficient $\lambda_{\mathrm{inc}}$.
The term is subtracted from the logits, so a positive $\lambda$ decreases attention as $\lvert p_i-p_j\rvert$ grows and a negative $\lambda$ increases it.

\Cref{tab:exp:bias_coef} reports the mean over the $8$ heads and $4$ layers of each stack ($32$ scalars per coordinate).
All branch means are positive except the Bias mode time mean ($-0.118$).
The typical head therefore penalizes a large \gls{toa} or incidence gap.
Every mode incidence scalar is positive in both models.
One Bias emitter head reaches $-0.832$ and instead raises attention as incidence differs, and some time heads in each stack do the same: those heads push the corresponding pulses together.
The trunk mean stays near zero.

RoPE-TOA+Bias keeps the same incidence pattern at smaller magnitude.
Mode time coefficients that were negative under Bias become positive once rotary encoding already carries $\Delta t$.

\begin{table}[H]
  \centering
  \small
  \caption{Learned bias coefficients after training. Each $\bar\lambda$ is the mean over the $8$ heads and $4$ layers of that stack. A positive coefficient decreases attention as $\lvert p_i-p_j\rvert$ grows (\cref{eq:method:attn_bias}).}
  \label{tab:exp:bias_coef}
  \begin{subtable}[t]{0.48\linewidth}
    \centering
    \caption{Bias}
    \begin{tabular}{@{}lcc@{}}
      \toprule
      Branch & $\bar\lambda_t$ & $\bar\lambda_{\mathrm{inc}}$ \\
      \midrule
      Trunk & $-0.026$ & {---} \\
      Mode & $-0.118$ & $0.877$ \\
      Emitter & $0.202$ & $0.229$ \\
      \bottomrule
    \end{tabular}
  \end{subtable}
  \hfill
  \begin{subtable}[t]{0.48\linewidth}
    \centering
    \caption{RoPE-TOA+Bias}
    \begin{tabular}{@{}lcc@{}}
      \toprule
      Branch & $\bar\lambda_t$ & $\bar\lambda_{\mathrm{inc}}$ \\
      \midrule
      Trunk & $0.037$ & {---} \\
      Mode & $0.107$ & $0.626$ \\
      Emitter & $0.167$ & $0.136$ \\
      \bottomrule
    \end{tabular}
  \end{subtable}
\end{table}
